\documentclass[12pt]{article}
\usepackage{amsmath,amssymb,amsfonts,amsthm}
\usepackage[margin=1in]{geometry}

\begin{document}

\begin{center}
{\Large\bfseries Chapter 9, Problem 27}\\[6pt]
Byron \& Fuller, \textit{Mathematics of Classical and Quantum Physics}
\end{center}

\bigskip

\textbf{Problem.}

(a) Show that if $H_n(x)$ is the $n$th Hermite polynomial satisfying Eq.~(5.105), then $\tilde{H}_n(x) = e^{-x^2/2}H_n(x)$ satisfies
\[
\left(-\frac{d^2}{dx^2} + x^2\right)\tilde{H}_n = (2n+1)\tilde{H}_n, \quad \text{and} \quad \int_{-\infty}^{\infty} \tilde{H}_n\,\tilde{H}_m\,dx = 0 \text{ for } n \neq m.
\]

(b) Show that $\{\tilde{H}_n\}$ are eigenfunctions of the Fourier transform $U_F$.

(c) Show $(U_F)^2 f(x) = f(-x)$, and hence the only possible eigenvalues of $U_F$ are $\pm 1, \pm i$.

(d) Characterize which Hermite functions have real vs.\ imaginary eigenvalues.

\bigskip
\textbf{Solution.}

\subsection*{(a) Harmonic oscillator eigenfunctions}

The Hermite polynomials satisfy the differential equation (Eq.~5.105):
\[
H_n'' - 2xH_n' + 2nH_n = 0.
\]

Define $\tilde{H}_n(x) = e^{-x^2/2}H_n(x)$. Compute:
\[
\tilde{H}_n' = e^{-x^2/2}(H_n' - xH_n),
\]
\[
\tilde{H}_n'' = e^{-x^2/2}(H_n'' - 2xH_n' + x^2 H_n - H_n).
\]

Therefore:
\[
-\tilde{H}_n'' + x^2\tilde{H}_n = e^{-x^2/2}(-H_n'' + 2xH_n' - x^2H_n + H_n + x^2H_n) = e^{-x^2/2}(-H_n'' + 2xH_n' + H_n).
\]

Using the Hermite equation $H_n'' - 2xH_n' = -2nH_n$:
\[
-\tilde{H}_n'' + x^2\tilde{H}_n = e^{-x^2/2}(2nH_n + H_n) = (2n+1)\tilde{H}_n.
\]

Therefore:
\[
\boxed{\left(-\frac{d^2}{dx^2} + x^2\right)\tilde{H}_n = (2n+1)\,\tilde{H}_n}.
\]

\textbf{Orthogonality:} Since $-d^2/dx^2 + x^2$ is a self-adjoint operator on $L_2(\mathbb{R})$, eigenfunctions with distinct eigenvalues are orthogonal:
\[
\int_{-\infty}^{\infty} \tilde{H}_n(x)\,\tilde{H}_m(x)\,dx = 0 \quad \text{for } n \neq m.
\]

(This also follows from the known orthogonality $\int H_n H_m e^{-x^2}\,dx = 0$ for $n \neq m$.) \hfill $\square$

\subsection*{(b) Eigenfunctions of the Fourier transform}

Denote the Fourier transform $(U_F \tilde{H}_n)(k) = \frac{1}{\sqrt{2\pi}}\int_{-\infty}^{\infty} \tilde{H}_n(x)\,e^{-ikx}\,dx$.

\textbf{Key insight:} The operator $\hat{O} = -d^2/dx^2 + x^2$ commutes with $U_F$. To see this:

Under the Fourier transform, $x \to i\frac{d}{dk}$ and $\frac{d}{dx} \to ik$. Therefore:
\[
U_F \hat{O} U_F^{-1} = -(ik)^2 + \left(i\frac{d}{dk}\right)^2 = k^2 - \frac{d^2}{dk^2} = \hat{O}.
\]

So $\hat{O}$ commutes with $U_F$. Since $\tilde{H}_n$ is an eigenfunction of $\hat{O}$ with eigenvalue $2n+1$, the function $U_F\tilde{H}_n$ is also an eigenfunction of $\hat{O}$ with the same eigenvalue.

If the eigenvalue $2n+1$ is non-degenerate (which it is for the harmonic oscillator---each level is non-degenerate), then $U_F\tilde{H}_n$ must be proportional to $\tilde{H}_n$:
\[
U_F\tilde{H}_n = c_n\,\tilde{H}_n
\]
for some constant $c_n$. Thus $\tilde{H}_n$ is an eigenfunction of $U_F$. \hfill $\square$

\subsection*{(c) $(U_F)^2 f(x) = f(-x)$ and eigenvalues}

Apply $U_F$ twice:
\begin{align}
(U_F^2 f)(x) &= \frac{1}{\sqrt{2\pi}}\int e^{-ikx}\left[\frac{1}{\sqrt{2\pi}}\int e^{-iky}f(y)\,dy\right]dk \\
&= \frac{1}{2\pi}\int f(y)\left[\int e^{-ik(x+y)}\,dk\right]dy = \int f(y)\,\delta(x+y)\,dy = f(-x).
\end{align}

So $U_F^2$ is the parity operator: $(U_F^2 f)(x) = f(-x)$.

Applying again: $(U_F^4 f)(x) = (U_F^2 f(-\cdot))(x) = f(-(-x)) = f(x)$, so $U_F^4 = I$.

If $U_F \tilde{H}_n = c_n \tilde{H}_n$, then $c_n^4 = 1$, giving:
\[
\boxed{c_n \in \{1, -1, i, -i\}}.
\]

\subsection*{(d) Classification by parity}

Since $U_F^2 \tilde{H}_n = c_n^2 \tilde{H}_n$ and $U_F^2 f(x) = f(-x)$:
\[
c_n^2 \tilde{H}_n(x) = \tilde{H}_n(-x) = (-1)^n \tilde{H}_n(x),
\]
using the parity property $H_n(-x) = (-1)^n H_n(x)$ (Hermite polynomials have definite parity).

Therefore $c_n^2 = (-1)^n$, giving:
\begin{itemize}
\item \textbf{$n$ even:} $c_n^2 = 1$, so $c_n = \pm 1$ (\textbf{real} eigenvalues).
\item \textbf{$n$ odd:} $c_n^2 = -1$, so $c_n = \pm i$ (\textbf{imaginary} eigenvalues).
\end{itemize}

More specifically, by direct computation or by tracking phases: $c_n = (-i)^n$, i.e.:
\[
c_0 = 1, \quad c_1 = -i, \quad c_2 = -1, \quad c_3 = i, \quad c_4 = 1, \quad \ldots
\]

The pattern repeats with period 4. In summary:
\begin{itemize}
\item $n \equiv 0 \pmod{4}$: eigenvalue $+1$.
\item $n \equiv 1 \pmod{4}$: eigenvalue $-i$.
\item $n \equiv 2 \pmod{4}$: eigenvalue $-1$.
\item $n \equiv 3 \pmod{4}$: eigenvalue $+i$.
\end{itemize}

\hfill $\blacksquare$

\end{document}
